(1) 令 $y(x,t)=X(x)\cos(\omega t+\phi)$。空间方程 $X''+k^2X=0$、$k=\omega/c$。左端固定使 $X=A\sin(kx)$，右端 Robin 条件得
\[
Tk\cos(kL)+K\sin(kL)=0.
\]
设 $\xi=kL,\ \beta=KL/T$，则正根满足 $\xi\cos\xi+\beta\sin\xi=0$，在 $\sin\xi\ne0$ 的根上等价于 $\tan\xi=-\xi/\beta$。

(2) 乘空间方程 $-TX''=\mu\omega^2X$ 以 $X$ 并在 $0$ 到 $L$ 积分，分部积分且 $X(0)=0,\ TX'(L)=-KX(L)$，得
\[
\mu\omega^2\int_0^L X^2dx=T\int_0^L(X')^2dx+KX(L)^2>0
\]
对任意非零 $X$ 成立，故 $\omega^2>0$。对两个不同本征值的方程交叉相乘并相减，边界项在 Robin 条件下抵消，于是 $(\omega_j^2-\omega_i^2)\int_0^L X_iX_j\,dx=0$，证明不同频率模正交。

(3) $K\to0$ 时右端自由，$X'(L)=0$，故 $k_nL\to(n+\tfrac12)\pi$。$K\to\infty$ 时有限能量要求 $X(L)=0$，右端固定，故 $k_nL\to(n+1)\pi$。两者不能简单按同一式直接令 $\beta=0$ 后忽略孤立零根；原始边界和正频率必须一起检查。

(4) 在 $\beta=0$ 正根附近设 $\xi=\xi_n^{(0)}+\delta$，$\xi_n^{(0)}=(n+\tfrac12)\pi$。线性化 $\xi\cos\xi+\beta\sin\xi=0$，利用 $\sin\xi_n^{(0)}=(-1)^n$ 得 $\delta=\beta/\xi_n^{(0)}+O(\beta^2)$，因此
\[
k_n=\frac1L\left[(n+\tfrac12)\pi+\frac{\beta}{(n+\tfrac12)\pi}+O(\beta^2)\right],\quad \omega_n=ck_n.
\]
弹簧增强约束使所有频率较自由端情形上移，与能量法一致。
